\documentclass[12pt,a4paper]{article}

\usepackage[utf8]{inputenc}
\usepackage{fontspec}
\usepackage{libertinus-otf}
\usepackage[french]{babel}
\usepackage{microtype}
\usepackage{icomma}

\setlength{\parindent}{0pt}

\usepackage[
    top=1.6cm,
    bottom=1.6cm,
    left=1.6cm,
    right=1.6cm,
    headheight=15pt,
    headsep=12pt,
    footskip=22pt
]{geometry}

\usepackage{fancyhdr}
\usepackage{lastpage}
\usepackage{tabularx}
\usepackage{tcolorbox}
\usepackage{enumitem}
\usepackage{eso-pic}
\usepackage{tikz}
\usepackage{xcolor}

\usepackage{amsmath, amssymb, esvect, mathtools}
\newcommand{\R}{\mathbb{R}}
\newcommand{\N}{\mathbb{N}}
\newcommand{\e}{\mathrm{e}}
\newcommand{\vect}[1]{\vv{#1}}

\setlength{\abovedisplayskip}{6pt}
\setlength{\belowdisplayskip}{6pt}
\setlength{\abovedisplayshortskip}{3pt}
\setlength{\belowdisplayshortskip}{3pt}

\newcolumntype{Y}{>{\centering\arraybackslash}X}

\newcommand{\exercice}[3][5]{%
    \noindent{\Large\textbf{Exercice #2}} \hfill \textbf{#1 points}\\[0.20cm]%
    {\normalsize\textit{#3}}%
    \par\vspace{0.28cm}\hrule\vspace{0.55cm}%
}
\newcommand{\partie}[1]{\vspace{0.45cm}\noindent\textbf{\large Partie #1}\par\vspace{0.25cm}}

\setlist[enumerate,1]{label=\textbf{\arabic*.}, leftmargin=*, itemsep=0.35em, parsep=0pt, topsep=0.2em}
\setlist[enumerate,2]{label=\textbf{\alph*.}, leftmargin=*, itemsep=0.25em, parsep=0pt, topsep=0.15em}

\pagestyle{fancy}
\fancyhf{}
\fancyhead[L]{\small\textit{Bac Blanc -- Spécialité Mathématiques}}
\fancyhead[C]{\small\textit{Sujet : Bienaymé-Tchebychev & Équations Différentielles}}
\fancyhead[R]{\small\textbf{Term\textsuperscript{spé}}}
\fancyfoot[C]{\thepage\ / \pageref{LastPage}}
\renewcommand{\headrulewidth}{0.4pt}
\renewcommand{\footrulewidth}{0.4pt}

\begin{document}

\thispagestyle{empty}

\AddToShipoutPictureFG*{%
    \begin{tikzpicture}[remember picture, overlay]
        \node[
            rotate=32,
            scale=1.1,
            text=black,
            opacity=0.06,
            font=\fontsize{72pt}{72pt}\bfseries\sffamily
        ] at (current page.center) {SP{\'E}CIALIT{\'E}};
    \end{tikzpicture}%
}

\noindent
\begin{minipage}[t]{0.52\textwidth}
    \textbf{\Large LYCÉE JEAN ROSTAND}\\[0.12cm]
    \textsc{Année scolaire 2026--2027}\\
    \textit{Classe de Terminale Générale}
\end{minipage}
\hfill
\begin{minipage}[t]{0.45\textwidth}
    \raggedleft
    \textbf{\large BACCALAURÉAT BLANC}\\[0.10cm]
    \textbf{Épreuve de spécialité}\\[0.10cm]
    \textit{Session de Printemps}\\[0.12cm]
    \small Date : Avril 2027
\end{minipage}

\vspace{0.30cm}
\hrule height 0.8pt
\vspace{0.45cm}

\vfill

\begin{center}
    {\LARGE \textbf{BACCALAURÉAT BLANC}}\\[0.20cm]
    {\large \textbf{SESSION 2027}}\\[0.15cm]
    {\textbf{\Large SUJET UNIQUE --- SPÉCIALITÉ MATHÉMATIQUES}}\\[0.45cm]

    \begin{tcolorbox}[colback=white, colframe=black, arc=0mm, boxrule=1.5pt, width=0.96\textwidth, top=4mm, bottom=4mm]
        \centering
        {\fontsize{26pt}{30pt}\selectfont\textbf{MATHÉMATIQUES}}\\[0.30cm]
        {\Large \textsc{Épreuve d'Enseignement de Spécialité}}
    \end{tcolorbox}

    \vspace{0.50cm}
    {\large \textbf{Durée de l'épreuve : 4 heures}}\\[0.25cm]
    \textbf{L'usage de la calculatrice avec \underline{mode examen actif} est autorisé.}\\[0.10cm]
    \textit{L'usage de la calculatrice sans mémoire, « type collège », est autorisé.}
\end{center}

\vfill

\begin{center}
\renewcommand{\arraystretch}{1.5}
\begin{tabularx}{0.96\textwidth}{|>{\centering\arraybackslash\bfseries}m{2.6cm}|*{4}{Y|}>{\centering\arraybackslash\bfseries}m{2.2cm}|}
    \hline
    Exercice & \textbf{Exercice 1} & \textbf{Exercice 2} & \textbf{Exercice 3} & \textbf{Exercice 4} & \textbf{Total} \\ \hline
    Thème    & Probabilités & Analyse & Nombres complexes & Analyse & --- \\ \hline
    Barème   & 4 points & 5 points & 6 points & 5 points & 20 points \\ \hline
\end{tabularx}
\end{center}

\vfill

\begin{center}
\begin{tcolorbox}[colback=gray!5, colframe=black, boxrule=1pt, arc=1mm, width=0.96\textwidth, top=3mm, bottom=3mm, left=4mm, right=4mm]
    \centering
    \textbf{\textsc{Consigne importante aux surveillants}}\\[0.15cm]
    \normalsize
    Nous remercions les surveillants de veiller à ce que chaque candidat inscrive lisiblement sur sa copie le \textbf{nom de son professeur de spécialité}.
\end{tcolorbox}
\end{center}

\vfill

\begin{tcolorbox}[colback=white, colframe=black, arc=0mm, boxrule=0.8pt, title=\textbf{\textsc{Consignes aux candidats}}, coltitle=black, colbacktitle=gray!15, top=3mm, bottom=3mm, left=4mm, right=4mm]
    \small
    \begin{itemize}[leftmargin=1.5em, itemsep=0.25em, label=---]
        \item Le sujet comporte \textbf{\pageref{LastPage} pages} numérotées de 1 à \pageref{LastPage}.
        \item Le candidat doit traiter les \textbf{4 exercices}.
        \item La qualité de la rédaction, la clarté et la précision des raisonnements seront prises en compte dans l'appréciation des copies.
        \item Tous les résultats doivent être justifiés, sauf mention contraire explicite.
        \item L'usage de tout document, téléphone ou appareil électronique connecté autre que la calculatrice autorisée est strictement interdit.
    \end{itemize}
\end{tcolorbox}

\newpage

\exercice[4]{1}{Thème : Probabilités -- Inégalité de Bienaymé-Tchebychev}

Soit $X$ une variable aléatoire suivant une loi binomiale $\mathcal{B}(n, p)$ avec $n \in \mathbb{N}^*$ et $p = 0,4$. On pose $Y_n = \frac{X}{n}$ la fréquence observée du succès sur $n$ épreuves indépendantes.

\vspace{0.45cm}

\begin{enumerate}
    \item Rappeler l'espérance $E(X)$ et la variance $V(X)$ de la variable aléatoire $X$. En déduire l'espérance $E(Y_n)$ et la variance $V(Y_n)$.
    
    \vspace{0.75cm}
    \item Énoncer l'inégalité de Bienaymé-Tchebychev pour une variable aléatoire $Z$ admettant une espérance et une variance finies.
    
    \vspace{0.75cm}
    \item En appliquant l'inégalité de Bienaymé-Tchebychev à la variable $Y_n$, démontrer que pour tout $\varepsilon > 0$ :
    \[
        P\left(\left|Y_n - 0,4\right| \geq \varepsilon\right) \leq \frac{0,24}{n \varepsilon^2}\,.
    \]
    
    \vspace{0.75cm}
    \item Déterminer le nombre minimal d'épreuves $n$ à réaliser pour garantir, avec une probabilité supérieure ou égale à $0,95$, que la fréquence observée $Y_n$ s'écarte de la probabilité théorique $0,4$ de moins de $0,05$.
\end{enumerate}

\newpage

\exercice[5]{2}{Thème : Analyse -- Équations différentielles et Intégration}

On considère l'équation différentielle $(E) : y' + 2y = 4\mathrm{e}^{-2x}$.

\vspace{0.45cm}

\begin{enumerate}
    \item Résoudre sur $\mathbb{R}$ l'équation homogène associée $(E_0) : y' + 2y = 0$.
    
    \vspace{0.55cm}
    \item Déterminer une solution particulière de $(E)$ de la forme $f_p(x) = ax\mathrm{e}^{-2x}$ où $a$ est un nombre réel.
    
    \vspace{0.55cm}
    \item En déduire l'ensemble de toutes les solutions de l'équation $(E)$.
    
    \vspace{0.55cm}
    \item Déterminer la solution unique $f$ de $(E)$ vérifiant la condition initiale $f(0) = 3$.
    
    \vspace{0.55cm}
    \item On considère la fonction $f$ définie sur $\mathbb{R}$ par $f(x) = (3 + 4x)\mathrm{e}^{-2x}$.
    \begin{enumerate}
        \item Calculer la limite de $f(x)$ en $+\infty$. Interpréter graphiquement ce résultat.
        \item Étudier les variations de $f$ sur $\mathbb{R}$ et dresser son tableau de variations complet.
        \item Calculer la valeur exacte de l'intégrale $I = \int_{0}^{1} f(x)\,\mathrm{d}x$.
    \end{enumerate}
\end{enumerate}

\newpage

\exercice[6]{3}{Thème : Nombres complexes et Géométrie}

Le plan complexe est muni d'un repère orthonormé direct $(O; \vec{u}, \vec{v})$.

\vspace{0.45cm}

\partie{A}

\begin{enumerate}
    \item Résoudre dans l'ensemble des nombres complexes l'équation :
    \[
        z^2 - 4z + 8 = 0\,.
    \]
    
    \item On considère les points $A$, $B$ et $C$ d'affixes respectives $z_A = 2 + 2\mathrm{i}$, $z_B = 2 - 2\mathrm{i}$ et $z_C = -2 + 2\mathrm{i}$.
    \begin{enumerate}
        \item Calculer le module et un argument de $z_A$ et de $z_C$.
        \item Placer ces points dans le repère et démontrer que le triangle $OAB$ est isocèle en $O$.
    \end{enumerate}
\end{enumerate}

\partie{B}

\begin{enumerate}
    \item Soit $f$ la transformation du plan qui à tout point $M$ d'affixe $z$ associe le point $M'$ d'affixe $z' = \mathrm{i}z + 2$.
    \begin{enumerate}
        \item Déterminer la nature et les éléments caractéristiques de la transformation $f$.
        \item Déterminer l'affixe du point $A'$ image de $A$ par $f$.
    \end{enumerate}
    
    \item On définit une suite de nombres complexes $(z_n)$ par $z_0 = 1$ et, pour tout entier naturel $n$, $z_{n+1} = \frac{1}{2}\mathrm{i}z_n + 1$.
    \begin{enumerate}
        \item Calculer $z_1$ et $z_2$.
        \item On pose, pour tout entier naturel $n$, $u_n = z_n - \frac{2}{2 + \mathrm{i}}$. Démontrer que $(u_n)$ est une suite géométrique dont on précisera le premier terme et la raison complexe.
        \item Exprimer $u_n$ puis $z_n$ en fonction de $n$, et étudier la convergence de la suite $(z_n)$.
    \end{enumerate}
\end{enumerate}

\newpage

\exercice[5]{4}{Thème : Analyse -- Étude de fonction et Intégration par parties}

Soit $g$ la fonction définie sur $]_0; +\infty[$ par $g(x) = x\ln(x) - x$.

\vspace{0.45cm}

\partie{A}

\begin{enumerate}
    \item Déterminer la limite de $g(x)$ en $0$ (par valeurs supérieures).
    
    \vspace{0.45cm}
    \item Étudier les variations de $g$ sur $]_0; +\infty[$ et dresser son tableau de variations.
    
    \vspace{0.45cm}
    \item Calculer l'intégrale $J = \int_{1}^{\mathrm{e}} \ln(x)\,\mathrm{d}x$ à l'aide d'une intégration par parties.
\end{enumerate}

\partie{B}

\begin{enumerate}
    \item Soit $h$ la fonction définie sur $]_0; +\infty[$ par $h(x) = x(\ln(x))^2 - 2x\ln(x) + 2x$.
    \begin{enumerate}
        \item Démontrer que $h$ est une primitive de la fonction $g$ sur $]_0; +\infty[$.
        \item Calculer l'aire exacte, exprimée en unités d'aire, du domaine plan délimité par la courbe représentative de $g$, l'axe des abscisses et les droites d'équations $x = 1$ et $x = \mathrm{e}$.
    \end{enumerate}
\end{enumerate}

\end{document}
